\subsection{包含 B 的 G 的子群}

对于任意子集 \(X\)，它是 \(S\) 的子集，我们将 \(W_X\) 记作 \(W\) 中由 \(X\) 生成的子群，并将 \(G_X\) 记作 \(BW_X B\) 中的双陪集 \(C(w)\)（其中 \(w \in W_X\)）的并集。我们有 \(G_\varnothing = B\) 且 \(G_S = G\)。

定理 3. a) 对于任意子集 \(X\)（\(S\) 的子集），集合 \(G_X\) 是 \(G\) 的一个子群，由 \(\bigcup_{s \in X} C(s)\) 生成。

\begin{enumerate}
\def\labelenumi{\alph{enumi})}
\setcounter{enumi}{1}
\item
  映射 \(\mathrm{X} \mapsto \mathrm{G}_{\mathrm{X}}\) 是从 \(\mathfrak{P}(\mathrm{S})\) 到 \(\mathrm{G}\) 的包含 \(\mathrm{B}\) 的子群集合的一个双射。
\item
  设 \((\mathrm{X}_i)_{i\in I}\) 是 X 的子集族。若 \(X = \bigcap_{i\in I}X_i\)，则 \(G_X = \bigcap_{i\in I}G_{X_i}\)。
\item
  设 X 和 Y 是 S 的两个子集。则 \(G_{X} \subset G_{Y}\)（分别为 \(G_{X} = G_{Y}\)）当且仅当 \(X \subset Y\)（分别为 X = Y）。
\end{enumerate}

显然 \(G_{X} = (G_{X})^{-1}\)；第 1 号的引理 1 表明 \(G_{X}.G_{X} \subset G_{X}\)；因此，考虑到定理 2 的推论 1，\(a)\) 得证。

映射 \(\mathrm{X} \mapsto \mathrm{G}_{\mathrm{X}}\) 的单射性由映射 \(\mathrm{X} \mapsto \mathrm{W}_{\mathrm{X}}\) 的单射性推出（参见~§1，第 8 号，定理 2）。反过来，设 \(\mathrm{H}\) 是 \(\mathrm{G}\) 的一个包含 \(\mathrm{B}\) 的子群。令 \(\mathrm{U}\) 为 \(w \in \mathrm{W}\) 且 \(C(w) \subset \mathrm{H}\) 的集合。我们有 \(\mathrm{H} = \mathrm{BUB}\)，因为 \(\mathrm{H}\) 是双陪集的并。令 \(\mathrm{X} = \mathrm{U} \cap \mathrm{S}\)；我们证明 \(\mathrm{H} = \mathrm{G}_{\mathrm{X}}\)。显然，\(\mathrm{G}_{\mathrm{X}} \subset \mathrm{H}\)。另一方面，令 \(u \in \mathrm{U}\)，并令 \((s_1, \ldots, s_q)\) 是 \(u\) 的一个约化分解。定理 2 的推论 3 蕴含 \(C(s_j) \subset \mathrm{H}\)，从而 \(s_j \in \mathrm{X}\)，其中 \(1 \leqslant j \leqslant q\)。因此，\(u \in \mathrm{W}_{\mathrm{X}}\)，且由于 \(\mathrm{H}\) 是各个 \(C(u)\) 的并（其中 \(u \in \mathrm{U}\)），我们有 \(\mathrm{H} \subset \mathrm{G}_{\mathrm{X}}\)，这就证明了 \(b\)。

断言 \(c)\) 和 \(d)\) 由 \(W_X\) 的相应性质推出（\(\S 1\)，第 8 号，定理 2）。

推论。集合 \(S\) 由这样的元素 \(w \in W\) 组成：\(w \neq 1\) 且 \(B \cup C(w)\) 是 \(G\) 的一个子群。

满足 \(w \in W\) 且 \(B \cup C(w)\) 是 G 的子群的元素，正是存在 \(X \subset S\) 使得 \(W_{X} = \{1, w\}\) 的那些元素。此外，若 \(w \neq 1\)，则必有 \(\text{Card}(X) = 1\)，即 \(w \in S\)。

备注。1) 上述推论表明，S 由 (G, B, N) 决定；因此，我们有时也允许自己说 (G, B, N) 是一个 Tits 系统，或 (B, N) 是 G 中的一个 Tits 系统。

命题 1. 设 X 是 S 的一个子集，N' 是 N 的一个子群，其在 W 中的像等于 W\_X。则，(G\_X, B, N', X) 是一个 Tits 系统。

我们有 \(G_{X} = BW_{X}B = BN'B\)，这表明 \(G_{X}\) 由 \(B \cup N'\) 生成。现在验证公理 (T1) 至 (T4) 是直接的。

命题 2. 设 \(\mathrm{X}, \mathrm{Y} \subset \mathrm{S}\) 且 \(w \in \mathrm{W}\)。我们有

\[
\mathrm{G} _ {\mathrm{X}} w \mathrm{G} _ {\mathrm{Y}} = \mathrm{BW} _ {\mathrm{X}} w \mathrm{W} _ {\mathrm{Y}} \mathrm{B}.
\]

令 \(s_1, \ldots, s_q \in X\) 且 \(t_1, \ldots, t_q \in Y\)。引理 1 表明

\[
\mathrm{C} (s _ {1} \dots s _ {q}). \mathrm{C} (w). \mathrm{C} (t _ {1} \dots t _ {q}) \subset \mathrm{BW} _ {\mathrm{X}} w \mathrm{W} _ {\mathrm{Y}} \mathrm{B},
\]

从而

\[
\mathrm{G} _ {\mathrm{X}} w \mathrm{G} _ {\mathrm{Y}} \subset \mathrm{BW} _ {\mathrm{X}} w \mathrm{W} _ {\mathrm{Y}} \mathrm{B}.
\]

反向包含关系显然成立。

备注。2) 记 \(G_X \backslash G / G_Y\) 为 \(G\) 中形如 \(G_X gG_Y\)（其中 \(g \in G\)）的子集的集合；并类似地定义 \(W_X \backslash W / W_Y\)。前一命题表明，典范双射 \(w \mapsto C(w)\)，从 \(W\) 到 \(B \backslash G / B\)，通过取商定义了一个双射 \(W_X \backslash W / W_Y \to G_X \backslash G / G_Y\)。

命题 3. 设 \(\mathbf{X} \subset \mathbf{S}\) 且 \(g \in \mathbf{G}\)。关系 \(g\mathrm{B}g^{-1} \subset \mathbf{G}_{\mathbf{X}}\) 蕴含 \(g \in \mathbf{G}_{\mathbf{X}}\)。

设 \(w \in W\) 使得 \(g \in C(w)\)。由于 B 是 \(G_X\) 的一个子群，假设 \(gBg^{-1} \subset G_X\) 蕴含 \(C(w).C(w^{-1}) \subset G_X\)，进而由定理 2 的推论 3 得 \(C(w) \subset G_X\)，所以 g 属于 \(G_X\)。

